\documentclass[11pt]{article}

\usepackage[margin=1.1in]{geometry}
\usepackage{amsmath,amssymb,amsthm}
\usepackage{booktabs}
\usepackage{enumitem}
\usepackage[colorlinks=true,linkcolor=blue!50!black,citecolor=blue!50!black,urlcolor=blue!50!black]{hyperref}

\newcommand{\N}{N}
\newcommand{\err}{\mathrm{err}}
\newcommand{\Unif}{\mathcal{U}}
\newcommand{\Strat}{\mathcal{S}}

\newtheorem{theorem}{Theorem}
\newtheorem{lemma}[theorem]{Lemma}
\newtheorem{remark}[theorem]{Remark}
\theoremstyle{definition}
\newtheorem{definition}[theorem]{Definition}

\title{Stratified Query Sampling for\\ Merkle-Committed Proximity Tests}
\author{Flock}
\date{August 5, 2026 --- revised October 15, 2026}

\begin{document}
\maketitle

\begin{abstract}
The query phase of a proximity test draws $q$ positions of a
Merkle-committed codeword, and its soundness analysis charges a cheating
prover $(1-\gamma)^q$: the probability that $q$ independent uniform draws
all avoid a $\gamma$-fraction of disagreeing positions. Inside a
recursive verifier, uniform sampling has a hidden cost: the commitment
node at which a query's authentication path terminates is data-dependent.
We show this cost is an artifact of the distribution, not of soundness.
Writing $q = \sum_i 2^{c_i}$ in binary and drawing, for each summand,
one uniform query inside each of the $2^{c_i}$ depth-$c_i$ subtrees, the
worst-case all-miss probability is \emph{exactly} $(1-\gamma)^q$. Every
path then terminates at a protocol-constant node, paths shorten from
$d$ to $d - c_i$ hashes, and the commitment shrinks to the
$2^{\lfloor \log_2 q \rfloor}$ nodes at the top set bit. Equal-size
strata are necessary, the binary decomposition is hash-minimal, and the
stratification survives FRI-style folding.
\end{abstract}

\section{Setting}

A Merkle tree of depth $d$ commits $\N = 2^d$ leaves; the node at depth
$k$ and index $j$ commits the leaf interval
$S^{(k)}_j = [\, j 2^{d-k}, (j+1) 2^{d-k})$. A \emph{depth-$c$ cap}
commitment publishes the $2^c$ nodes at depth $c$
\cite{plonky2}; an opening of leaf $x$ carries $d - c$ siblings and is
checked against cap word $\lfloor x / 2^{d-c} \rfloor$.

Proximity tests of the Ligero and FRI families \cite{ligero,ligerito,fri}
share a soundness template for their query phase. For every cheating
prover and every transcript prefix up to the moment the queries are
drawn, the analysis exhibits a set $A \subseteq [\N]$ of density at most
$1-\gamma$ (the positions where the committed oracle agrees with what the
verifier will check) such that the verifier accepts only if all $q$
queries land in $A$. The set $A$ may depend on everything before the
queries, including the verifier's earlier randomness, but not on the
query randomness; nothing is assumed about its structure. The phase's
error is therefore governed by one quantity of the sampling distribution.

\begin{definition}
For a distribution $\mathcal{D}$ on $[\N]^q$ and $\gamma \in (0,1]$, let
\[
\err_\gamma(\mathcal{D}) \;=\;
\max_{|A| \le (1-\gamma)\N}\;
\Pr_{(x_1,\dots,x_q) \leftarrow \mathcal{D}}[\, x_1, \dots, x_q \in A \,].
\]
\end{definition}

For $q$ independent uniform draws, $\err_\gamma(\Unif^q) = \mu(A)^q$ at
$|A| = \lfloor (1-\gamma)\N \rfloor$, i.e.\ $(1-\gamma)^q$ up to
integrality; this is the benchmark. We write $\mu(A) = |A|/\N$ and
$\mu_S(A) = |A \cap S| / |S|$.

The cost of uniformity appears when the verifier is replayed in a
circuit. A uniform query authenticates against a cap word whose index is
unknown at circuit-build time, so the circuit must bind the path's
terminal dynamically: a multiplexer over the cap, per-query terminals as
public inputs, or hashing the cap to a root in-circuit and extending
every path to full depth. Each pays for the data dependence. We ask for
distributions $\mathcal{D}$ with $\err_\gamma(\mathcal{D}) \le
(1-\gamma)^q$ whose draws are confined to protocol-constant subtrees.

\section{The construction}

\begin{lemma}[equal strata]\label{lem:amgm}
Fix $c \le d$ and $r \ge 1$, and let $\mathcal{D}_{c,r}$ draw,
independently, $r$ uniform positions in each of the $2^c$ depth-$c$
subtrees $S_1, \dots, S_{2^c}$. For every $A$,
\[
\Pr_{\mathcal{D}_{c,r}}[\text{all in } A]
= \prod_{j} \mu_{S_j}(A)^{r}
\le \Big( \frac{1}{2^c} \sum_{j} \mu_{S_j}(A) \Big)^{r 2^c}
= \mu(A)^{r 2^c},
\]
with equality iff $\mu_{S_j}(A) = \mu(A)$ for all $j$.
\end{lemma}

\begin{proof}
The $r 2^c$ draws are independent, and a draw that is uniform on $S_j$
lands in $A$ with probability $\mu_{S_j}(A)$; the $r$ draws in stratum
$j$ contribute the factor $\mu_{S_j}(A)^r$, which gives the first
equality. For the inequality, apply the AM--GM inequality to the $2^c$
numbers $\mu_{S_1}(A), \dots, \mu_{S_{2^c}}(A) \in [0,1]$:
\[
\prod_{j} \mu_{S_j}(A) \;\le\; \Big( \frac{1}{2^c} \sum_j \mu_{S_j}(A) \Big)^{2^c},
\]
and raise both sides to the power $r$. For the last equality, the strata
partition $[\N]$ into $2^c$ parts of equal size $\N / 2^c$, so
\[
\frac{1}{2^c} \sum_j \mu_{S_j}(A)
= \frac{1}{2^c} \sum_j \frac{|A \cap S_j|}{\N/2^c}
= \frac{|A|}{\N} = \mu(A).
\]
AM--GM is an equality iff all $2^c$ numbers coincide, and since their
mean is $\mu(A)$ this means $\mu_{S_j}(A) = \mu(A)$ for every $j$.
\end{proof}

\begin{definition}[stratified schedule]
A schedule is a multiset of depths $c_1 \ge \dots \ge c_t$ in $[0,d]$
with query count $q = \sum_i 2^{c_i}$; its distribution $\Strat$ draws,
independently across summands, one uniform position in each depth-$c_i$
subtree (a repeated depth is a summand with $r = 2$). The
\emph{canonical} schedule for $q \le \N$ is the binary representation of
$q$, with $c_1 = \lfloor \log_2 q \rfloor$.
\end{definition}

Query $(i,j)$ is confined to the protocol-constant stratum
$S^{(c_i)}_j$ and consumes $d - c_i$ random bits.

\begin{theorem}[exactness]\label{thm:main}
For every schedule with query count $q$,
$\err_\gamma(\Strat) \le (1-\gamma)^q$, with equality whenever
$(1-\gamma)2^{d-c_1}$ is an integer.
\end{theorem}

\begin{proof}
Fix $A$. By independence across summands,
$\Pr_{\Strat}[\text{all in } A] = \prod_i P_i(A)$ with
$P_i(A) = \prod_j \mu_{S^{(c_i)}_j}(A)$, and Lemma~\ref{lem:amgm} gives
$P_i(A) \le \mu(A)^{2^{c_i}} \le (1-\gamma)^{2^{c_i}}$; multiply.

A priori $\max_A \prod_i P_i(A) \le \prod_i \max_A P_i(A)$ could be
strict, since one $A$ faces every summand. It is not: each $P_i$ is
maximized by an $A$ of density $1-\gamma$ on every depth-$c_i$ stratum,
and an $A$ uniform at the finest granularity (the depth-$c_1$ strata) is
uniform on every coarser stratum, as those are unions of fine ones. Such
an $A$ exists exactly under the integrality condition and attains every
factor's maximum at once.
\end{proof}

Hence stratified sampling is a drop-in replacement for independent
uniform sampling under the template: no query count, security parameter,
or query-for-grinding substitution moves, and a schedule exists for
every integer $q$.

\begin{remark}[equal strata are necessary]
The lemma fails for unequal strata. Two strata of relative size $9/10$
and $1/10$, one query each, $\mu(A) = 9/10$: the adversary takes
$\mu_{S_2}(A) = 1$, $\mu_{S_1}(A) = 8/9$ and all-miss probability
$8/9 > 0.81$. The weighted AM--GM
$\prod_j \mu_{S_j}(A)^{w_j} \le \mu(A)$ shows the general rule: queries
per stratum proportional to stratum size (Neyman's proportional
allocation \cite{neyman}, here for a worst-case tail rather than a
variance). One query per stratum forces equal sizes, hence dyadic
subtrees and a decomposition of $q$ into powers of two.
\end{remark}

\begin{remark}[no free lunch; fewer collisions]
Stratification does not beat $(1-\gamma)^q$ (equality case of
Theorem~\ref{thm:main}); sampling without replacement does, by a
lower-order factor the usual accounting ignores. Stratification does
reduce collisions: queries of one summand lie in disjoint strata, so at
least $2^{c_1}$ positions are always distinct.
\end{remark}

\begin{remark}[the schedule is part of the statement]
The verifier must fix the schedule from $q$, never from the proof. A
schedule that leaves one depth-$c$ stratum unqueried lets the adversary
hide a $2^{-c}$-fraction of disagreement there, lowering the effective
disagreement to $(\gamma - 2^{-c})/(1-2^{-c})$ and voiding the test when
$2^{-c} \ge \gamma$. The invariant is that every summand covers the
leaves with equal strata, each queried the same number of times.
\end{remark}

\section{Payoff in a recursive verifier}

Take the canonical schedule $c_1 > \dots > c_t$. The commitment is the
$2^{c_1}$ cap nodes. A query of summand $i$ hashes $d - c_i$ rows and
terminates on a wire holding node $(c_i, j)$: for $i = 1$ a cap word,
otherwise a node of a shared \emph{upper tree} that folds the cap
upward from depth $c_1$ to depth $c_t$ in $2^{c_1} - 2^{c_t}$ rows
(every layer is either a terminal layer at a set bit or an operand of
the next; no root is formed unless $q$ is odd). Siblings above the cap
are verifier-recomputable witness, so the proof carries $d - c_1$
siblings per query for every summand. Every terminal is a copy
constraint to a fixed wire; the only data-dependent part of a query is
its $d - c_i$ low bits driving the path's left/right selection. Total
hash rows are
\[
\textstyle\sum_i 2^{c_i}(d - c_i) \;+\; 2^{c_1} - 2^{c_t}.
\]
The binary decomposition minimizes this over all schedules with the
same $q \le \N$: merging two depth-$c$ summands into one depth-$(c+1)$
summand saves $2^{c+1}$ path rows and adds at most $2^c$ upper-tree
rows (when $c$ was the top depth). The counter-pressure, a larger
absorbed cap, is second order since absorbed words cost a fraction of a
row each.

At a representative level, $d = 19$, $q = 90 = 64 + 16 + 8 + 2$:

\begin{center}
\begin{tabular}{lrrr}
\toprule
& root & cap $+$ in-circuit root & stratified \\
\midrule
absorbed commitment (words) & 1 & 128 & 64 \\
circuit path rows & $1{,}710$ & $1{,}710$ & $1{,}236$ \\
circuit tree rows & 0 & 127 & 62 \\
\textbf{circuit hash rows} & $\mathbf{1{,}710}$ & $\mathbf{1{,}837}$ & $\mathbf{1{,}298}$ \\
proof siblings per level & $1{,}710$ & $1{,}080$ & $1{,}170$ \\
query randomness (bits) & $1{,}710$ & $1{,}710$ & $1{,}236$ \\
\bottomrule
\end{tabular}
\end{center}

Stratified is the first design that beats a root commitment in-circuit
($-24\%$) while keeping the cap's proof-size advantage; against the cap
design it saves $29\%$ of the query-phase rows and halves the absorbed
commitment.

\begin{remark}[query reuse under folding]
FRI-style protocols reuse a query across levels, $x \mapsto x \bmod
\N/2$. In bit-reversed storage (leaf $p$ holds index $\mathrm{rev}_d(p)$,
so folded pairs are adjacent) the leaf position evolves as $p \mapsto
\lfloor p/2 \rfloor$, which preserves the top $c$ bits of $p$. A query
confined to depth-$c$ stratum $j$ stays in stratum $j$ at every later
level until the tree depth drops below $c$, so one schedule, chosen at
the first level, is inherited by the rest. For $a$-ary trees the same
analysis holds with the base-$a$ representation of $q$.
\end{remark}

\begin{thebibliography}{9}\small

\bibitem{ligero}
S.~Ames, C.~Hazay, Y.~Ishai, M.~Venkitasubramaniam.
\emph{Ligero: lightweight sublinear arguments without a trusted setup}.
CCS 2017.

\bibitem{ligerito}
A.~Novakovi\'c, G.~Angeris. \emph{Ligerito}. [TODO: confirm venue.]

\bibitem{fri}
E.~Ben-Sasson, I.~Bentov, Y.~Horesh, M.~Riabzev.
\emph{Fast Reed--Solomon interactive oracle proofs of proximity}.
ICALP 2018.

\bibitem{plonky2}
Polygon Zero. \emph{Plonky2: fast recursive arguments with PLONK and
FRI}. 2022. [TODO: confirm the cap-commitment reference.]

\bibitem{neyman}
J.~Neyman. \emph{On the two different aspects of the representative
method}. J.\ Royal Statistical Society 97(4), 1934.

\end{thebibliography}

\end{document}
